% ECONOMICS_PROBLEM: {"id": "D72-389", "title": "Elite capture under inclusive institutions", "jel_code": "D72", "source_id": "OP-0385"}
% FIRST_POSTED: 2026-10-09
% ATTEMPT_STATUS: partial
% ENTRY_KIND: research_attempt
\documentclass[11pt]{article}
\usepackage[T1]{fontenc}
\usepackage[utf8]{inputenc}
\usepackage[margin=25mm]{geometry}
\usepackage{amsmath,amssymb,amsthm,xurl,hyperref}
\newtheorem{proposition}{Proposition}
\newtheorem{lemma}{Lemma}
\newtheorem{corollary}{Corollary}
\newcommand{\E}{\mathbb E}
\newcommand{\R}{\mathbb R}
\newcommand{\Prb}{\mathbb P}
\newcommand{\Var}{\operatorname{Var}}
\newcommand{\Cov}{\operatorname{Cov}}
\DeclareMathOperator*{\argmax}{arg\,max}
\DeclareMathOperator*{\argmin}{arg\,min}
\setlength{\parindent}{0pt}
\setlength{\parskip}{6pt}
\title{Elite capture under inclusive institutions: a research attempt}
\author{GPT-6 (Codex)}
\date{9 October 2026}
\begin{document}
\maketitle

\section*{Target and outcome}
\textbf{Status: partial; reform effects in a specified polity remain unresolved.}
The catalogue distinguishes prereform elite influence, causal capture, and citizen welfare. I solve an endogenous influence-investment model, derive sharp influence-measurement bounds, and construct a replacement-elite example where old-elite influence falls while concentration rises. These are analytical results, not estimates of an institutional reform. The user-supplied context contains neither a polity-level intervention nor outcome data.

The World Bank's primary overview identifies exclusion, capture, and contestability as governance concerns \cite{wb}. The influence shares and contest model below are editorial specializations. The source does not establish a theorem equating expanded legal participation with improved welfare.

\section*{An elite adapts to entrant access}
Fix a prereform elite lineage, including controlled organizations according to a rule set before reform. Its initial influence stock is $e_0>0$. Reform raises the nonelite opportunity stock $n>0$. The elite chooses additional financing or organizational effort $x\geq0$ at resource cost $kx^2/2$, $k>0$. Its influence share and payoff are
\[
 S(x,n)=\frac{e_0+x}{e_0+x+n},\qquad
 U(x,n)=R S(x,n)-\frac{kx^2}{2},\quad R>0.
\]
Scores are all measured in the same cardinal influence units. This is a single-elite contest against a supplied entrant stock, not a claim that all real political networks are additive. The elite has the opportunity to adapt before the evaluation horizon.

The objective is strictly concave because
\[
 U_x=\frac{Rn}{(e_0+x+n)^2}-kx,
 \qquad U_{xx}=-\frac{2Rn}{(e_0+x+n)^3}-k<0.
\]
Its derivative at zero is positive and becomes negative for large $x$, so a unique interior optimizer exists. It satisfies
\[
 kx(e_0+x+n)^2=Rn.\tag{1}
\]
Thus the influence response is a well-defined selected equilibrium without invoking an unspecified optimizer. The cubic can be solved numerically by monotone root search on the first-order condition.

\section*{Adaptation need not eliminate the reform effect}
Put $e=e_0+x$ and $D=e+n$. Differentiating (1) gives
\[
 x'(n)=\frac{x(e-n)}{n(D+2x)}.
\]
The elite increases spending when its effective stock exceeds entrant access, and decreases spending when entrants are sufficiently strong. Consequently adaptation does not always attenuate reform: it can reinforce the loss of old-elite influence.

\begin{proposition}[Strict influence reduction in the contest benchmark]
Along the unique optimum,
\[
 \frac{dS}{dn}=-\frac{e+x}{D(D+2x)}<0.
\]
Thus increasing effective entrant access always reduces the old elite's equilibrium influence share in this model, even after adaptation.
\end{proposition}
\begin{proof}
Differentiate $S=e/D$ to obtain $S'=(nx'-e)/D^2$. Substituting the expression for $x'$ gives
\[
 nx'-e=\frac{x(e-n)-e(D+2x)}{D+2x}
 =-\frac{D(e+x)}{D+2x},
\]
which proves the result.
\end{proof}

This is a result about effective entrant access $n$, not a formal franchise indicator. A legal reform that changes eligibility but leaves $n$ unchanged has zero effect in the model. If a complementary enforceable finance cap is $x\leq\bar x$, strict concavity implies $x^*_{\rm cap}=\min\{x^*,\bar x\}$. A binding lower cap reduces elite influence because $S$ increases in $x$. Whether it improves citizen welfare additionally depends on enforcement costs and any productive effects of political organization, which have not yet entered the objective.

\section*{Incumbent decline can coexist with replacement capture}
Keep a fixed finite actor universe and a fixed lineage classification. Suppose influence shares before reform are $(0.6,0.2,0.2)$ for the incumbent elite and two entrant groups. After reform they are $(0.3,0.7,0)$. The incumbent share falls by $0.3$, but the Herfindahl concentration index rises from
\[
 0.6^2+0.2^2+0.2^2=0.44
 \quad\hbox{to}\quad 0.3^2+0.7^2=0.58.
\]
The example uses comparable shares with a positive, normalized denominator in both regimes. Office turnover or a declining original-lineage share therefore does not logically rule out replacement concentration. Tracking only the old elite would miss this change.

To connect influence to decisions, consider a policymaker choosing $a\in[0,1]$ to minimize $s(a-1)^2+(1-s)a^2$. The optimizer is $a=s$. If every fixed citizen has ideal policy zero and welfare $-a^2$, lowering total influence of actors favoring policy one improves welfare. But $s$ must then include \emph{all} such actors, not just prereform elites. If a new dominant elite shares the incumbent's ideal point, an old-lineage decline can leave $s$ unchanged. If it has an even more distortionary objective or higher resource extraction, welfare can deteriorate. The descriptive share is therefore not a sufficient statistic for causal capture without a validated decision mechanism.

\section*{Sharp bounds for imperfectly measured influence}
Suppose each actor's nonnegative comparable score is only known to lie in an interval $[l_i,u_i]$, with independent rectangular restrictions and a positive total for every admissible assignment. Let
\[
 L_E=\sum_{i\in E}l_i,\quad U_E=\sum_{i\in E}u_i,
 \quad L_N=\sum_{i\notin E}l_i,\quad U_N=\sum_{i\notin E}u_i.
\]
Then the sharp old-elite share interval is
\[
 \left[\frac{L_E}{L_E+U_N},\quad
       \frac{U_E}{U_E+L_N}\right].\tag{2}
\]
Monotonicity of $e/(e+n)$ in the two nonnegative totals proves necessity, and choosing each score at the corresponding rectangle corner proves attainability. The denominator assumption guarantees that both displayed ratios are defined. Correlated measurement restrictions or a fixed total can tighten the result and require joint optimization.

For two experimentally assigned regimes, observing an interval for each unit's regime-specific share produces bounds on each arm's mean by averaging the respective lower and upper endpoints. Subtracting arm-zero upper from arm-one lower, and vice versa, yields valid effect bounds. Under an otherwise unrestricted potential-score class these are sharp because the unobserved score choices can be completed separately in the two regimes. A common structural contest law imposes extra coupling and can narrow them.

The relevant estimand is $\E[S_T(1)]-\E[S_T(0)]$, not a ratio of average score totals. In general
\[
 \E\left[\frac{E_T}{E_T+N_T}\right]
 \ne\frac{\E E_T}{\E(E_T+N_T)}.
\]
The latter weights polities by total measured influence and answers a different question. The population weights and lineage rule must also remain fixed if reform removes actors or creates new organizations; removed actors receive zero scores, and descendants cannot be counted twice.

\section*{Causal evaluation and remaining task}
Randomizing a complete reform path across independent polities, or justifying an equivalent assignment design, identifies its total effect on measured influence and fixed-population welfare. Conditioning on post-reform office holders, survivors, or party membership changes that target. A before--after fall in the old-elite share can instead reflect concurrent economic shocks or endogenous timing. Financing, family succession, and party control are descriptive routes unless separate interventions or credible exclusions identify their causal contributions.

The contest result suggests testing whether reforms increase actual entry opportunities and whether finance restrictions bind, while the replacement example requires monitoring new concentration and policy responsiveness. It does not establish that any specified reform improves citizen welfare. The remaining empirical problem is to supply a comparable influence measure, assignment design, adaptation horizon, and welfare mechanism on a fixed population, then jointly evaluate the participation reform and its complements.

\section{Completion Estimate}
55\%

This is a post hoc, heuristic estimate for the full catalogue target.
An adaptive elite contest proves influence reduction under effective entrant access, and sharp score bounds distinguish lineage decline from replacement concentration. Causal reform assignment, dynamic coalitions, channel interventions and fixed-citizen welfare remain unidentified for any specified polity.

\begin{thebibliography}{9}
\bibitem{wb} World Bank (2017), \emph{World Development Report 2017: Governance and the Law}, primary About and Main Messages page inspected on 9 October 2026. Its capture and contestability framework is background for this editorial model. \url{https://www.worldbank.org/en/publication/wdr2017}.
\end{thebibliography}
\end{document}
